\section{Introduction}\label{sec:intro}

Finding the cheapest way from one place to another is one of the most common
problems in computing. A navigation system planning a drive from Dhaka to
Sylhet, a game character walking around a river and a router forwarding a
packet across a network all solve it, often on maps with millions of places and
with an answer needed in a fraction of a second. Two classical algorithms
approach the problem from opposite directions. Dijkstra's algorithm always finds
a cheapest route but searches evenly in every direction, whereas greedy
best-first search heads straight for the destination but may settle for a
longer route.

The A* search algorithm of Hart, Nilsson and Raphael~\cite{hart1968formal}
combines the strengths of both and is the standard method for this problem when
an estimate of the remaining cost is available~\cite{russell2020aima}. It ranks
every place that it reaches by the cost of the best route found so far plus an
estimate of the cost that remains, such as the straight-line distance to the
destination. As long as this estimate never exceeds the true cost, A* returns a
cheapest route, yet it examines far fewer places than Dijkstra's algorithm. In
this report we derive A* from its two predecessors, prove when it is correct,
analyze its running time and memory, and explain why it examines fewer places.
We confirm the analysis on a small example, on a grid with an obstacle and on a
map of the 64 districts of Bangladesh.

\subsection{Pathfinding}\label{sec:pathfinding}

We model a map as a graph: places are vertices, allowed moves are edges, and
every edge has a nonnegative cost, such as a length or a travel time. A trip
begins at a \emph{source} and ends at a \emph{destination}.

\begin{definition}[Pathfinding problem]\label{def:pathfinding}
Given a graph whose edges have nonnegative costs, a source vertex $s$ and a
destination vertex $t$, the \emph{pathfinding problem} asks for a path from $s$
to $t$ whose total cost is minimum.
\end{definition}

Figure~\ref{fig:pathfinding} shows two instances. In the road map, the cheapest
route from $A$ to $B$ costs 8 although it visits the most places. In the
lattice, a building blocks the direct route, and the cheapest route climbs over
the roof in 12 moves instead of eight.

\begin{figure}[htbp]
  \centering
  \begin{subfigure}[b]{0.40\textwidth}\centering
    \begin{tikzpicture}[x=1cm, y=1cm]
      \coordinate (A) at (0,0.9);   \coordinate (p) at (1.6,1.8);
      \coordinate (q) at (1.6,0);   \coordinate (r) at (3.2,0.9);
      \coordinate (B) at (4.8,0.9);
      \draw[algastar!45, line width=5pt, line cap=round, line join=round]
        (A) -- (p) -- (q) -- (r) -- (B);
      \foreach \a/\b/\w in {A/p/2, A/q/4, p/q/1, p/r/5, q/r/2, r/B/3}
        {\draw[storyedge] (\a) -- node[storyweight] {\w} (\b);}
      \node[storystart] at (A) {$A$};
      \node[storynode]  at (p) {$p$};
      \node[storynode]  at (q) {$q$};
      \node[storynode]  at (r) {$r$};
      \node[storygoal]  at (B) {$B$};
    \end{tikzpicture}\par\vspace{6pt}
    \caption{A road map. The cheapest route, $(A, p, q, r, B)$,
      costs 8.}\label{fig:pf-graph}
  \end{subfigure}\hfill
  \begin{subfigure}[b]{0.56\textwidth}\centering
    \begin{tikzpicture}[cam={16}{28}{1.02}{1.02}, line cap=round, line join=round]
  \fill[black!3] (-0.5,-0.5,0) -- (5.5,-0.5,0) -- (5.5,3.5,0) -- (-0.5,3.5,0) -- cycle;
  \foreach \y in {0,...,3}{\foreach \z in {0,1,2}{\draw[latedge] (0,\y,\z) -- (1,\y,\z);}}
  \foreach \x in {0,1}{%
    \foreach \z in {0,1,2}{\draw[latedge] (\x,0,\z) -- (\x,3,\z);}
    \foreach \y in {0,...,3}{\draw[latedge] (\x,\y,0) -- (\x,\y,2);}}
  \foreach \x in {0,1}{\foreach \y in {0,...,3}{\foreach \z in {0,1,2}{%
    \node[latnode] at (\x,\y,\z) {};}}}
  \fill[gwall!78!white] (1.55,-0.45,0) -- (3.45,-0.45,0) -- (3.45,-0.45,1.45) -- (1.55,-0.45,1.45) -- cycle;
  \fill[gwall]          (3.45,-0.45,0) -- (3.45,3.45,0)  -- (3.45,3.45,1.45)  -- (3.45,-0.45,1.45) -- cycle;
  \fill[gwall!45!white] (1.55,-0.45,1.45) -- (3.45,-0.45,1.45) -- (3.45,3.45,1.45) -- (1.55,3.45,1.45) -- cycle;
  \foreach \wx in {1.8,2.3,2.8}{\foreach \wz in {0.3,0.85}{%
    \fill[white!55!gwall] (\wx,-0.45,\wz) -- (\wx+0.3,-0.45,\wz)
      -- (\wx+0.3,-0.45,\wz+0.3) -- (\wx,-0.45,\wz+0.3) -- cycle;}}
  \foreach \y in {0,...,3}{\draw[latedge] (1,\y,2) -- (4,\y,2);}
  \foreach \x in {2,3}{\draw[latedge] (\x,0,2) -- (\x,3,2);}
  \foreach \y in {0,...,3}{\foreach \z in {0,1,2}{\draw[latedge] (4,\y,\z) -- (5,\y,\z);}}
  \foreach \x in {4,5}{%
    \foreach \z in {0,1,2}{\draw[latedge] (\x,0,\z) -- (\x,3,\z);}
    \foreach \y in {0,...,3}{\draw[latedge] (\x,\y,0) -- (\x,\y,2);}}
  \foreach \x in {2,3}{\foreach \y in {0,...,3}{\node[latnode] at (\x,\y,2) {};}}
  \foreach \x in {4,5}{\foreach \y in {0,...,3}{\foreach \z in {0,1,2}{%
    \node[latnode] at (\x,\y,\z) {};}}}
  \draw[edgpath]
    (0,0,0) -- (0,1,0) -- (1,1,0) -- (1,1,1) -- (1,1,2) -- (2,1,2)
    -- (3,1,2) -- (4,1,2) -- (4,1,1) -- (4,1,0) -- (4,2,0) -- (5,2,0)
    -- (5,3,0);
  \foreach \p in {(0,1,0),(1,1,0),(1,1,1),(1,1,2),(2,1,2),(3,1,2),
                  (4,1,2),(4,1,1),(4,1,0),(4,2,0),(5,2,0)}{%
    \node[latnode, fill=cbGreen, draw=white] at \p {};}
  \node[sball] at (0,0,0) {};
  \node[tag, text=gstart, below left=2pt] at (0,0,0) {\textbf{A}};
  \node[gball] at (5,3,0) {};
  \node[tag, text=ggoal, right=4pt] at (5,3,0) {\textbf{B}};
  \begin{scope}[shift={(-0.9,-1.3,0)}]
    \draw[-{Stealth[length=1.4mm]}, neutralg, line width=0.5pt] (0,0,0) -- (0.6,0,0)
      node[tag, text=neutralg, right=-1pt] {$x$};
    \draw[-{Stealth[length=1.4mm]}, neutralg, line width=0.5pt] (0,0,0) -- (0,0.7,0)
      node[tag, text=neutralg, above right=-2pt] {$y$};
    \draw[-{Stealth[length=1.4mm]}, neutralg, line width=0.5pt] (0,0,0) -- (0,0,0.6)
      node[tag, text=neutralg, above=-1pt] {$z$};
  \end{scope}
\end{tikzpicture}

    \caption{A three-dimensional lattice. A building removes the points inside
      it, and the cheapest route climbs over the roof.}\label{fig:pf-3d}
  \end{subfigure}
  \caption{Two instances of the pathfinding problem.}\label{fig:pathfinding}
\end{figure}

\subsection{Applications}\label{sec:applications}

Table~\ref{tab:applications} lists typical applications of pathfinding.

\begin{table}[H]
  \centering
  \caption{Applications of pathfinding.}\label{tab:applications}
  \small
  \renewcommand{\arraystretch}{1.25}%
  \begin{tabular}{@{}P{2.5cm}P{2.6cm}P{3.1cm}P{5.6cm}@{}}
    \toprule
    Application & Vertices & Cost of an edge & Example\\
    \midrule
    Road navigation & road junctions & length or travel time of a road
      & planning the fastest drive from Dhaka to Sylhet\\
    Video games & tiles of the map & difficulty of the terrain entered
      & a character that walks around a river to a bridge instead of wading
        through it\\
    Computer networks & routers & delay or load of a link
            & the OSPF protocol\tablefootnote{OSPF (Open Shortest Path First) is the
        routing protocol with which the routers inside one network exchange
        link costs and compute the shortest path to every destination.}, in
        which every router runs Dijkstra's algorithm to choose the paths of data
        packets\\
    Robotics & positions in free space & distance or energy
      & warehouse robots that carry shelves to packing stations without
        colliding\\
    Molecular biology & shapes of a protein & energy needed to change shape
      & the Folding@home project, which searches networks of protein shapes for
        likely folding pathways\\
    Puzzles & board positions & one per move
      & the fewest moves that solve the 15-puzzle or Rubik's
        Cube~\cite{korf1985, korf1997}\\
    \bottomrule
  \end{tabular}
\end{table}

\subsection{The need for pathfinding algorithms}\label{sec:need}

Listing every route and keeping the cheapest is infeasible, because the number
of routes grows exponentially with the size of the map. On a $50 \times 50$
grid there are about $2.5 \times 10^{28}$ routes between opposite corners that
only move right or up; checking a billion per second would take about 800
billion years. Pathfinding algorithms instead grow routes outward from the
source and discard a route as soon as a cheaper route to the same place is
known.

\subsection{Requirements of a pathfinding algorithm}\label{sec:requirements}

A pathfinding algorithm must serve every pair of source and destination, so we
judge it by two requirements.
\begin{description}[font=\normalfont\bfseries, leftmargin=1.5em, labelindent=0pt]
  \item[Optimality.] Whenever the destination can be reached from the source,
    the algorithm returns a cheapest route between them.
  \item[Efficiency.] The algorithm finds this route while examining as little
    of the map as possible.
\end{description}
